\section{第41届复赛第3题：油水混合物的表面张力与微观表面层模型}
\subsection{题干}
题目给出石油与含表面活性物质的水混合后，表面张力系数与温度 $T$、压强 $p$ 的经验关系：
\[\frac{\sigma}{10^{-3}\,\mathrm{N\,m^{-1}}}=60+0.065\frac{T}{\mathrm K}-24.0\frac{p}{\mathrm{bar}}+3.15\frac{p^2}{\mathrm{bar}^2},\qquad \beta=7.28\times10^{-3}\,\mathrm{bar\,K^{-1}}.\]
石油初温为 $100^\circ\mathrm C$，密度、比热容分别为水的 $80\%$、$60\%$。恒容抽油并注入等体积水时，要求在表面张力最小时确定注入水温。第（2）问假定表面张力完全源于表面层和内部粒子邻近数的差异，表面层邻近数为内部的 $\zeta$ 倍，要求导出表面张力表达式。
\subsection{参考解答整理}
恒容条件给出 $dp=\beta\,dT$；又因参考状态有 $p=\beta T$，代入题设关系后得到
\[\frac{\sigma}{10^{-3}\,\mathrm{N\,m^{-1}}}=60-110\times10^{-3}\frac{T}{\mathrm K}+167\times10^{-6}\frac{T^2}{\mathrm K^2}.\]
令 $d\sigma/dT=0$，有 $T_f=329\,\mathrm K$，即 $t_f\approx56^\circ\mathrm C$。该二次关系的极小点及混合过程由下图辅助说明；该图为解析辅助图，并非原题图。
\begin{figure}[htbp]\centering\begin{tikzpicture}[>=Stealth,scale=.86]
  \draw[->,thick] (0,0)--(6.2,0) node[right] {$T/\mathrm{K}$};
  \draw[->,thick] (0,0)--(0,4.2) node[above] {$\sigma/(10^{-3}\,\mathrm{N\,m^{-1}})$};
  \draw[very thick,blue!70!black,domain=.35:5.75,samples=80] plot (\x,{.22*(\x-3.5)*(\x-3.5)+.7});
  \draw[densely dashed] (3.5,0)--(3.5,.7); \fill (3.5,.7) circle(1.5pt);
  \node[below] at (3.5,0) {$T_f=329\,\mathrm K$};
  \node[right,align=left] at (4.45,1.15) {$\sigma(T)$最小\\[-2pt]$t_f\approx56^\circ\mathrm C$};
  \draw[->,thick] (1.05,3.15)--(3.25,1.15) node[midway,above,sloped] {油与水混合、恒容};
  \node[align=center] at (1.05,3.65) {石油初温\\$t_i=100^\circ\mathrm C$};
  \node[below,align=center] at (3.5,-.72) {解析辅助图：由表面张力二次关系确定混合终温};
\end{tikzpicture}
\caption{第41届复赛第3题解析辅助图：表面张力极小值对应的混合终温。}\end{figure}
对等体积的石油和水作热量守恒，
\[c_i\rho_i(t_i-t_f)=c_w\rho_w(t_f-t_w).\]
代入 $c_i=0.6c_w$、$\rho_i=0.8\rho_w$、$t_i=100^\circ\mathrm C$、$t_f\approx56^\circ\mathrm C$，得到注入水温 \[\boxed{t_w\approx35^\circ\mathrm C}.\]

第（2）问中，设内部每个粒子有 $n$ 个邻近粒子，表面层有 $\zeta n$ 个。形成单位面积表面所需功对应表面层粒子的相互作用能减少；令特征粒子间距为 $D$，则 $D^2=\Delta A_S/N_S$。下图仅用于说明该微观模型的符号关系。
\begin{figure}[htbp]\centering\begin{tikzpicture}[>=Stealth,scale=.86]
  \fill[blue!7] (-3,-1.8) rectangle (3,0);
  \draw[thick] (-3,0)--(3,0); \node[right] at (3,0) {液面};
  \foreach \x/\y in {-2.5/-.45,-1.7/-.8,-.85/-.4,0/-.85,.9/-.42,1.75/-.85,2.55/-.42,-2.2/-1.4,-1.25/-1.25,-.35/-1.5,.55/-1.25,1.45/-1.5,2.35/-1.25} \fill[blue!55] (\x,\y) circle(.16);
  \foreach \x in {-2.4,-1.2,0,1.2,2.4} \fill[blue!55] (\x,.22) circle(.16);
  \draw[<->] (-2.4,.55)--(-1.2,.55) node[midway,above] {$D$};
  \draw[->] (2.05,.85)--(1.2,.27) node[above right,align=left] {表面层粒子\\邻近数为 $\zeta n$};
  \draw[->] (2.15,-2.5)--(1.35,-1.35) node[below right,align=left] {液体内部粒子\\邻近数为 $n$};
  \draw[<->] (-2.85,0)--(-2.85,-.58) node[midway,left] {$D$};
  \node[below,align=center] at (0,-2.85) {解析辅助图：表面层粒子数减少与特征间距 $D$};
\end{tikzpicture}
\caption{第41届复赛第3题解析辅助图：表面层与内部粒子的邻近关系。}\end{figure}
由 $L_m=N_A(|\epsilon|/2)n$、$D^3N_A=\mu/\rho$，最终获得
\[\boxed{\sigma=(1-\zeta)\frac{L_m}{N_A}\left(\frac{\rho N_A}{\mu}\right)^{2/3}}.\]
\subsection{转写依据}
题干、关系式与参数依据 \texttt{CPHO41-Q03\_transcription\_v1.md}；答案公式和数值依据答案第7至9页与 \texttt{CPHO41-Q03\_formula\_check\_v1.csv}。原题无独立题图；本条目的两图均明确标注为解析辅助图。
